% Figure method: generated by figures/rse_fig_method.py; placed where the text first discusses it.
\begin{figure}[!tbp]
  \centering
  \includegraphics[width=\textwidth]{../figures/out/rse_fig_method.pdf}
  \caption{\textbf{The hindcast audit, step by step, on one example.}
  The example is the Volga--Ural interfluve coast, the grid GEBCO\_2026 and the level bin centred at \metric{erl.hero.e3_b029.level_m}~m (bins are \metric{params.bin_width_m}~m wide; levels on the DAHITI reference).
  (a)~Left: daily basin level from satellite altimetry (grey line) and the \metric{erl.hero.e3_b029.n_acquisitions} Landsat acquisitions of \metric{erl.hero.e3_b029.first_year}--\metric{erl.hero.e3_b029.last_year} whose level falls in the bin (shaded band).
  The acquisitions are split by their order into two interleaved halves (filled circles, half A; open circles, half B).
  Right: water frequency, the share of the valid acquisitions in which a \metric{params.stage2_res_m}~m cell was classified as water (blank where fewer than \metric{params.min_valid_obs} acquisitions are valid).
  The black line is the observed shoreline: the edge of the cells whose frequency is at least \metric{params.contour_level} and that connect to the open sea.
  (b)~GEBCO\_2026 height minus the level, after the grid's vertical offset $c$ (\metric{h1.shape_offset.gebco2026_m}~m), which was fitted to the flooded area on training blocks.
  The black line is the predicted shoreline: the edge of the area flooded from the open sea at the level.
  A cell is flooded when a path from the open sea reaches it without crossing ground higher than the level.
  Hatched: cells below the level that the flood does not reach.
  (c)~Left: the predicted shoreline cells, coloured by their distance to the nearest observed shoreline cell; black, the observed shoreline; grey dashed and dotted lines, the shorelines observed by half A and by half B.
  Right: cumulative distributions, over the whole sector at this level, of the distances between predicted and observed shoreline cells (blue; distances from both shorelines pooled) and of the distances between the shorelines of the two halves (grey).
  The displacement $D$ is the median of the first set; the observation noise floor $Q$ is the \metric{params.err_pct_hi}th percentile of the second.
  (d)~$D/Q$ at each of the \metric{h1.n_evaluable.shape.gebco2026.interfluve} levels at which both $D$ and $Q$ can be measured in this sector (filled, $D > Q$; open, $D \le Q$; ringed, the example level of a--c).
  A grid fails the shoreline test in a sector when $D > Q$ at more than half of these levels; GEBCO\_2026 exceeds $Q$ at \metric{h1.n_exceeding.shape.gebco2026.interfluve} of them here.
  Scale bars \metric{params.rse_fig_method_scale_km}~km.}
  \label{fig:rse_method}
\end{figure}
